\section{Completion of the proof}\label{sec:completion}
It remains to evaluate the determinant comparison at the parameters of
Table~\ref{tab:parameters}. For the size estimate we use
\begin{equation}\label{eq:spectral-parameters}
 (a,v,\beta,a_0)=\left(\frac{39}{40},\frac{19}{50},-1,\frac1{200}\right),
 \qquad (N_t,N_z,N_m)=(128,256,512).
\end{equation}
The last two integers are the mesh sizes for the column-order capacity
bound and the intervals of possible $m/R$. The finite choices of affine
composition parameters are the same as in Appendix~C of~\cite{chung498}. These values satisfy the finite hypotheses of
Section~\ref{sec:comparison}. At each upper radial endpoint, the
calculation checks the moment and concavity conditions before using the
Schur bound.


For the arithmetic estimate we retain only the two- and three-factor
composite corrections. The cutoff for the derivative slope is
\[
 q_{\rm cap}=\min\{J-1,\,m_D^4-1,\,m_D\lfloor P/m_D\rfloor\},
 \qquad m_D=(D+1)/4.
\]
It is $46\,799\,995$ in the first parameter family, $47\,923\,198$ at
$D=379$, and $57\,289\,760$ at $D=347$. In each case the factor bounds
and column capacities in the arithmetic estimate are satisfied.

Table~\ref{tab:gaps} gives rigorous enclosures for the difference
$G=\arith-\metric$ between the arithmetic lower bound and the size upper
bound in Proposition~\ref{prop:comparison}. The calculations use directed
rounding, and the displayed endpoints are rounded outward.
\begin{table}[ht]
\centering
\begin{tabular}{rr}
\toprule
$D$ & Enclosure for $G$\\
\midrule
$347$ & $[0.00146083,\;0.00147313]$\\
$379$ & $[0.01186524,\;0.01187031]$\\
$419$ & $[0.00417435,\;0.00417636]$\\
$443$ & $[0.04790491,\;0.04790692]$\\
$467$ & $[0.11463345,\;0.11463546]$\\
$491$ & $[0.17810049,\;0.17810250]$\\
\bottomrule
\end{tabular}
\caption{The six positive comparison margins, with two- and three-factor
composite corrections.}
\label{tab:gaps}
\end{table}

\begin{proof}[Proof of Theorem~\ref{thm:main}]
Suppose $h(-D)=1$ and $D>331$. The bound $D\leq498$ from~\cite{chung498}
and the elementary reduction in Section~\ref{sec:intro} restrict $D$ to
$347,379,419,443,467,491$. Section~\ref{sec:primes} supplies the
auxiliary inert prime and verifies the hypotheses of
Theorem~\ref{thm:exact-rank}, so the interpolation matrix has a nonzero
maximal minor in each case. Proposition~\ref{prop:comparison} requires
$G\leq0$, whereas every lower endpoint in Table~\ref{tab:gaps} is positive.
This excludes all six possibilities.
\end{proof}

The bound $331$ shows that a sharper nonvanishing criterion allows the
existing determinant estimates to go further. The arithmetic estimate
still uses class number one to control non-inert primes by norm
representations. A general argument for Landau--Siegel zeros would
require a substitute for that input.
